% ============================================================
% UNIDAD 3: PFD Y MFD
% ============================================================
\clearpage
\section*{UNIDAD 3: EL PFD Y EL MFD}
\subsection*{Presentación de datos, alertas, banderas y modos reversionarios}

\begin{recuadro}
\textbf{Objetivo de la unidad:} interpretar la información que presentan el PFD y el MFD, sus
alertas y sus indicaciones de falla, y conocer los modos reversionarios y los instrumentos de
reserva, desde el punto de vista del técnico que debe verificarlos.
\end{recuadro}

% ------------------------------------------------------------
\section*{1. EL PFD}

\begin{center}
\begin{tikzpicture}[every node/.style={font=\scriptsize}]
  % marco
  \draw[very thick, fill=black!85] (0,0) rectangle (10,6.4);
  % horizonte
  \begin{scope}
    \clip (1.6,1.9) rectangle (8.4,6.2);
    \fill[blue!45] (0,4.05) rectangle (10,6.4);
    \fill[brown!60] (0,1.9) rectangle (10,4.05);
    \draw[white, thick] (0,4.05) -- (10,4.05);
    \foreach \y/\l in {4.55/5,5.05/10,5.55/15} { \draw[white] (4.6,\y) -- (5.4,\y); \node[white] at (5.75,\y) {\l}; }
    \foreach \y in {3.55,3.05} { \draw[white] (4.6,\y) -- (5.4,\y); }
  \end{scope}
  \draw[yellow, line width=2pt] (3.7,4.05) -- (4.6,4.05) -- (4.8,3.85);
  \draw[yellow, line width=2pt] (6.3,4.05) -- (5.4,4.05) -- (5.2,3.85);
  \fill[yellow] (5,4.05) circle (0.06);
  % cinta de velocidad
  \fill[gray!60] (0.3,2.1) rectangle (1.4,6.0);
  \foreach \y/\v in {2.6/90,3.35/100,4.05/110,4.75/120,5.45/130} { \node[white] at (0.8,\y) {\v}; }
  \fill[green!60!black] (1.3,2.1) rectangle (1.4,4.6);
  \draw[white, fill=black] (0.25,3.85) rectangle (1.45,4.25); \node[white] at (0.85,4.05) {\textbf{112}};
  % cinta de altitud
  \fill[gray!60] (8.6,2.1) rectangle (9.7,6.0);
  \foreach \y/\v in {2.6/4400,3.35/4450,4.75/4550,5.45/4600} { \node[white] at (9.15,\y) {\v}; }
  \draw[white, fill=black] (8.55,3.85) rectangle (9.75,4.25); \node[white] at (9.15,4.05) {\textbf{4500}};
  \draw[magenta, thick] (8.6,4.75) -- (8.75,4.75);
  % HSI
  \draw[white, thick] (5,0.75) circle (0.95);
  \foreach \a in {0,30,...,330} { \draw[white] (\a:0.95) ++(5,0.75) -- ++(\a:-0.15); }
  \node[white] at (5,1.55) {N};
  \draw[magenta, very thick, -{Stealth}] (5,0.05) -- (5,1.45);
  \draw[white, fill=black] (4.6,1.75) rectangle (5.4,2.05); \node[white] at (5,1.9) {360$^\circ$};
  % barra superior
  \fill[gray!30] (0,6.1) rectangle (10,6.4);
  \node at (2,6.25) {NAV1 / COM1};
  \node at (5,6.25) {Modos del piloto automático};
  \node at (8.2,6.25) {GPS / distancia};
  % etiquetas
  \node[anchor=east, align=right] at (-0.15,4.05) {Cinta de\\velocidad};
  \node[anchor=west, align=left] at (10.15,4.05) {Cinta de\\altitud};
  \node[anchor=west, align=left] at (10.15,2.7) {Variómetro\\(junto a la cinta)};
  \node[anchor=west, align=left] at (10.15,5.5) {Altitud\\seleccionada};
  \node[anchor=east, align=right] at (-0.15,0.75) {HSI con\\curso y CDI};
  \node[anchor=east, align=right] at (-0.15,5.5) {Horizonte\\artificial};
\end{tikzpicture}
\end{center}

\begin{itemize}
  \item \textbf{Indicador de actitud:} el horizonte ocupa toda la pantalla (a veces con
        \textbf{visión sintética} del terreno en 3D). Escala de cabeceo, arco de alabeo,
        indicador de deslizamiento y, opcionalmente, director de vuelo.
  \item \textbf{Cinta de velocidad:} a la izquierda, con la velocidad actual en un recuadro, los
        arcos de colores (blanco de flaps, verde normal, amarillo de precaución, raya roja de
        nunca exceder), tendencia de velocidad y velocidad verdadera.
  \item \textbf{Cinta de altitud:} a la derecha, con la altitud seleccionada, el ajuste
        barométrico y, junto a ella, la velocidad vertical.
  \item \textbf{HSI:} rosa de rumbos con el curso seleccionado, la desviación (CDI) del VOR, ILS o
        GPS, el indicador de rumbo seleccionado y la indicación de la fuente de navegación.
  \item \textbf{Barras de datos:} frecuencias de radio, modos del piloto automático, datos de GPS,
        temperatura exterior, código del transpondedor, hora.
\end{itemize}

% ------------------------------------------------------------
\section*{2. EL MFD}

Muestra lo que no entra en el PFD: \textbf{mapa móvil} con aeropuertos, espacios aéreos y la ruta;
\textbf{páginas de motor} (en aviones livianos, a menudo en una franja permanente); tránsito
(ADS-B o TCAS); meteorología (radar, datos por enlace satelital o ADS-B); terreno (TAWS);
\textbf{cartas de aproximación} y planos de aeropuertos; planificación de vuelo, listas de chequeo y
páginas de configuración y estado de los equipos. Muchas de estas funciones dependen de
\textbf{bases de datos} que deben estar vigentes.

% ------------------------------------------------------------
\section*{3. ALERTAS, AVISOS Y BANDERAS}

\begin{center}
\small
\begin{tabular}{|p{0.18\linewidth}|p{0.18\linewidth}|p{0.52\linewidth}|}
\hline
\textbf{Nivel} & \textbf{Color} & \textbf{Significado} \\ \hline
Advertencia (\textit{warning}) & Rojo & Requiere acción inmediata (por ejemplo, presión de aceite
baja, falla del alternador con baja tensión). Suele ir con aviso sonoro. \\ \hline
Precaución (\textit{caution}) & Ámbar / amarillo & Requiere atención (por ejemplo, combustible
bajo, falla de un sensor). \\ \hline
Aviso (\textit{advisory}) & Blanco o celeste & Información sobre el estado de un sistema. \\ \hline
\end{tabular}
\end{center}

\paragraph*{Banderas e indicación de dato no válido.} Cuando un dato no es confiable, el sistema
\textbf{no lo muestra}: lo reemplaza por una \textbf{cruz roja} (\textit{red X}) o una bandera
(``ATTITUDE FAIL'', ``AHRS ALIGN'', ``NO GPS''). Para el técnico, cada cruz roja es una pista: indica
qué fuente de datos se perdió (AHRS, ADC, magnetómetro, sensor de motor, comunicación con otra
unidad). El sistema también registra las fallas en sus páginas o registros de mantenimiento.

Otras alertas: velocidad de pérdida (con avisador de pérdida y, en algunos equipos, indicador de
ángulo de ataque), sobrevelocidad, alertas de terreno (TAWS/GPWS), tránsito, altitud y
desviación de los mínimos.

% ------------------------------------------------------------
\section*{4. MODOS REVERSIONARIOS E INSTRUMENTOS DE RESERVA}

\begin{itemize}
  \item \textbf{Modo reversionario:} si falla una pantalla, la otra pasa a mostrar la información
        esencial del PFD junto con el motor (automáticamente o con un botón rojo de reversión).
  \item \textbf{Redundancia de sensores:} los sistemas de mayor categoría duplican los AHRS y las
        ADC y permiten seleccionar la fuente.
  \item \textbf{Instrumentos de reserva:} con alimentación independiente y batería propia
        (habitualmente de 30 minutos a varias horas según el equipo y el requisito), deben permitir
        terminar el vuelo con seguridad.
  \item \textbf{Alimentación eléctrica:} las barras esenciales y las baterías de reserva de los
        equipos se verifican en el mantenimiento: si la batería interna de un instrumento de reserva
        no pasa su prueba de capacidad, el equipo no está en condiciones.
\end{itemize}

% ------------------------------------------------------------
\section*{5. ACTIVIDADES DE LA UNIDAD}
\begin{enumerate}
  \item En una imagen de un PFD, identificar y anotar cada elemento y su fuente de datos (ADC,
        AHRS, magnetómetro, GPS, radio de navegación).
  \item Para cada una de estas cruces rojas, indicar la probable unidad o sensor en falla:
        actitud, rumbo, velocidad y altitud a la vez, solo el rumbo, solo la EGT del cilindro 3.
  \item Con un simulador de vuelo con PFD y MFD (o el modo entrenamiento de un equipo), practicar
        el cambio de fuente de navegación, el ajuste barométrico y el modo reversionario.
  \item Buscar en el manual del piloto de un G1000 qué ocurre al presionar el botón de reversión.
\end{enumerate}
